\section{当前构建的修订验证}
\label{supp:major-revision}
生产代码快照 \texttt{9879f64} 保留了本地回收补丁
\texttt{bf71c4d}。新增测试夹具和网络驱动程序的新增部分
分别记录指纹；\texttt{721800c} 的测量数据保持
不变。工件目录 \texttt{major-revision-2026-10-03}
包含源码清单、原始日志、解析后的账本和审稿问题关闭记录。

\noindent\textbf{有限容量与 Intent 边界。}
确定性回归测试获取真实成员批准，通过
\texttt{FinalizeBlock}/\texttt{Commit} 执行区块，并在成员
跟随之前，利用保留的空后继区块认证执行结果。
测试框架固定执行顺序和区块时间。

在 $B_g=300$、每名成员有一个 Worker 的配置下，四个独立的
100 CAL 请求各获得两份批准，在没有形成 QC 的情况下，
耗尽全部四名成员各自的 200 CAL 容量。一次 150 CAL 的外部储备增加
在不删除锁的情况下，为每名成员增加 100 容量；随后同一组请求
得以完成，每名成员最终满足 Available$=300$ 和 Reserved$=0$。
全诚实成员下的同输入 2/2 分裂在相同增资之后仍然阻塞。

在用户付费和已获授权的组织付费两种费用模式下，
分别运行两轮跨组织 Intent 冲突。组织 A 的储备
从 300 降至 200，再降至 100 CAL；公开状态中的 Spent 达到 200，
Reserved 回到零。同样的三名源交易签署者各自在本地保留 200，
第四名成员仍保有未使用的容量；新请求无法从该切片
获得法定人数批准。经过认证的失败结果不会释放
被拒绝的批准。其原输入在公开账本中仍未被消费，
且在本地保持锁定。用户付费运行还保留两个完整的
10,000 FUEL 输入，这与各自 1,000 FUEL 的托管上限不同。

第一个后代输出为第二个 A 侧请求提供资金，最终
后代输出通过 B 成功消费。每个 B 侧赢家使用
不同的初始输入，测试框架在下一轮开始前完成补偿。
独立求和在每个已提交前缀检查 Open 缺口、证书覆盖、
授权项用量以及 CAL/FUEL 守恒。
重复补偿不会产生新的变更、效果或费用输出。
这些结果展示了有限的储备损失与保守的
认证停顿，而不是网络攻击速率。

\noindent\textbf{保留状态的重启与原始历史重放。}
使用真实门限适配的回归测试覆盖
补偿--表示--偿还以及
补偿--偿还--表示两种顺序。重放副本在
高度 2 停止，重新打开保留的持久化账本，并使用
已经包含修订区块体的 BlockStore 中的
\texttt{Original} 材料追赶。在每个固定命令历史内，
各重放高度的应用状态根和账本行均与参考结果一致。
这些测试覆盖本地持久状态重启和原始历史访问，
而不是网络下载，也不验证不同命令顺序下
应用状态根的相等性。

\noindent\textbf{费用与锁定资金。}
该轨迹中，每笔成功的普通付款支出 94 FUEL，
涉及补偿的消费方付款支出 99。两轮加上最终后代付款
共支出 480 FUEL，其中 430 用于奖励，50 被销毁；
托管上限合计为 5,000，返还 4,520。
用户付费轨迹另行保留未执行源交易的两个
10,000 FUEL 输入。组织付费轨迹
从 B 的已授权费用账户扣减 480，而 A 的公开 FUEL
余额不变。A 的本地费用预留仍然保留。CAL 损失、
费用支出和费用输入锁定是不同的量。

\noindent\textbf{公开入口与原子性覆盖。}
wire4 分派器接纳 DirectSubmission、储备增加、时钟
推进和受限修复命令。DirectSubmission 在执行共同的
消费状态检查之前，重建并验证当前组织的 QC、
所有者授权和直接输入证书。旧版普通转账
要求 CommitteeRoute，并被 wire4 分派器排除。
组织付费要求所用授权项的 Policy Subject 与
付款人匹配；未经批准的代付在成员路径和公开路径上均被拒绝。
零票／两票构造探针在规范编码阶段失败；
既有公开验证负向测试另行覆盖格式错误或
不匹配的证书。补偿和表示不能提供
任意消费所有者资金的替代入口。

同步批准边界测试在真实数据库提交之后、
控制权返回签署者之前立即中断。重新打开
同一存储后，预留和输入锁仍然保留；同一请求
可以在不再次扣减的情况下获得签名。既有测试在重启后拒绝
冲突请求，在跟随者应用失败时同时回滚游标与状态，
并确保跨重启的恢复仅生效一次。这些
检查保留已提交状态，而不是模拟其丢失。

\subsection{等待模式时间戳审计}
\noindent\textbf{跳间观测开销。}
驱动程序每 5\,ms 轮询一次，区块跟随者每 25\,ms 轮询一次；
250\,ms 是提交参数，而不是其中任何一个轮询间隔。
将原始 600 条付款记录与全部四个副本的
应用提交记录匹配后，可识别出每轮等待模式运行中的
99 次跳间等待。从接收到公开观测的时间区间分别累计为
56.851、58.921 和 57.031\,s。
对于相同的跳，从最早副本提交到驱动程序观测的
时间区间分别累计为 29.542、29.681 和 29.347\,s，
逐跳 P50 约为 297\,ms。后一时间区间涵盖后继区块头
认证、获取与验证、本地记录、观测以及
副本间差异；由于缺少钱包内部时间戳，
无法进一步拆分。从观测到下一笔付款构造的交接时间，
每轮分别仅累计 0.273、0.314 和 0.315\,ms。

因此，14.15 倍这一比值比较的是该配置下测得的
两种端到端模式。它既不是纯共识
加速比，也不是仅由轮询造成的开销；
这些时间区间也不能得出反事实的事件驱动结果。

\subsection{归档字节导出示例}
可运行的 \texttt{TestHistoricalExportConsumer} 生成本地导出，
内容包括前缀、坐标、完整 BlockID、Revision、Task ID、
原始及已获授权的区块体字节、历史资金来源和当前义务
状态。它验证 Task 所指定的精确区块体与修订版本、
原始部件集身份、所有者签名、QC 和输入 CH 关系。
迟到偿还会改变经济查询结果，但不会改变已获授权的字节。
这是已验证本地历史的一个示例消费方，并不声称它是
已部署客户端或远程认证协议。两条读取路径
都查询永久决定和义务。

\subsection{变色龙哈希实例化}
\label{sec:ch-params}

\noindent\textbf{参数。} 所用哈希为 RSA 变色龙哈希，模数 $N$ 为 2048 位，公钥指数为 $e=65537$。对于上下文 $c$ 和消息 $m$，摘要为 $\mathrm{CH}=H(c,m)\cdot r^{e}\bmod N$，其中随机值 $r$ 属于 $\mathbb{Z}_N^{*}$。给定 $(c,m,r)$ 的摘要和新消息 $m'$，适配后的随机值 $r'=r\,(H(c,m)/H(c,m'))^{d}\bmod N$ 满足 $H(c,m')\,r'^{e}=H(c,m)\,r^{e}$，其中 $d=e^{-1}$ 按群阶取模。由于正确生成的密钥满足 $\gcd(e,\varphi(N))=1$，映射 $r\mapsto r^e$ 是 $\mathbb{Z}_N^{*}$ 上的置换。因此，每个消息与摘要组合都确定唯一的有效开口。

\noindent\textbf{全域哈希。} $H$ 读取一个 SHAKE256 输出流，其输入依次为前缀 \texttt{UTXO\_CH\_RSA\_FDH\_V3}、256 字节的模数、$\mathrm{u64be}(|c|)$、上下文 $c$、$\mathrm{u64be}(|m|)$ 和消息 $m$。它从流中连续取出 256 字节的数据块，并返回第一个非零、小于 $N$ 且在模 $N$ 下可逆的值。长度前缀使上下文与消息的联合编码具有单射性。区块部件的上下文依赖高度和部件索引，而不依赖修订计数器。

\noindent\textbf{门限陷门。} 可信分发者使用 CIRCL \texttt{GenerateKey} 生成两个不同的安全素数，并检查 $e$ 在模 $\varphi(N)$ 下的逆元；随后使用 \texttt{Deal}，通过 Cloudflare CIRCL v1.6.5 的 \texttt{tss/rsa}（\texttt{Deal(4,3,true)}），将 $\widetilde d=e^{-1}\bmod\mu$ 拆分为四份，其中 $\mu=(p-1)(q-1)/4$。实验从不持久化完整私钥。每名份额持有者将公开比值 $H(c,m)/H(c,m')$ 作为原始群元素进行签名，而不是通过带填充的 RSA 签名 API。\texttt{CombineSignShares} 重新组合份额，组合者至多尝试四个三成员组，直到最终开口通过公开关系式验证。该构造遵循 Shoup~\cite{shoup} 的门限 RSA 方法。原始群元素接口通过其最终公开关系式进行检查；仅引用该库并不构成该接口或账本授权的证明。

\noindent\textbf{原始群元素接口的正确性。} 记 $u=H(c,m)/H(c,m')$，$\Delta=4!=24$。对诚实份额进行插值得到 $W$，满足 $W^e=u^{4\Delta^2}$。由于 $\gcd(4\Delta^2,e)=\gcd(2304,65537)=1$，选取整数 $\alpha,\beta$，使其满足 $\alpha4\Delta^2+\beta e=1$。组合者得到 $z=W^\alpha u^\beta$，满足 $z^e=u$，因此 $r'=rz$ 保持承诺不变。模 $\mu$ 下的共享指数 $\widetilde d$ 不同于上文用于表达唯一根的全群指数 $d$。这说明了对可逆元群中比值进行诚实组合的正确性，而不是对可公开组合的原始比值作出不可伪造性主张。

\noindent\textbf{假设与范围。} Next 的公钥构造器只检查模数是否为 2048 位奇数，以及是否满足 $e=65537$。它不证明模数构造正确，因此论证要求分发者正确生成并拆分密钥。比值是公开的，且可通过代数运算组合，因此变色龙哈希本身不能证明某次替换已获得对应授权。适配只有通过已提交的补偿决定和精确的 Task 才能生效。这些对象的检查独立于哈希。原始区块体由修订版本~0、开口~1 的原始部件，以及正文历史读取接口所述的完整身份检查共同绑定。这种绑定不依赖陷门的保密性。共识安全仍然依赖基础协议的 BFT、认证和状态守卫。

\subsection{表示构造与批次检查}
构造器读取最新已提交的逻辑表示，不依赖物理安装，并选择槽位仍保存 \texttt{OriginalFunding} 的已决定输出；决定在提交时已经确保资金足以支付，因此构造器不模拟任何经济步骤。对于每个条目，它根据网络和已消费输出导出储备扣款身份，并将槽位中的 \texttt{OriginalFunding} 替换为带有开口的储备引用，将同一交易的各槽位修改应用到同一副本中。序列化长度和所有非目标字节保持不变，包括此前已获授权的替换。输入适配逐槽位执行；合并后的区块体所涉及的每个区块部件仅适配一次，未改变的部件则保留原有承诺与证明。

份额端点绑定至固定的委员会身份，每名持有者从自身已提交的决定中导出所请求的替换，并拒绝没有相应决定的槽位。收集者通过实际密码学组合测试数量有界的三成员候选组，且仅在每个受影响部件均通过验证时接受候选组，因此格式错误或无效响应会淘汰相应候选组，而其他成员提供的三份有效贡献仍可使用。组合成功后取消尚未完成的请求。开口提供承诺兼容性，已提交决定提供授权。

\begin{figure}[t]
  \small
  \begin{algorithmic}[1]
    \REQUIRE 批次 $B$、状态 $S$、执行高度
    \STATE $b\gets\operatorname{BatchID}(B)$；如果已存在成功的 Task $b$，
    返回无新增效果
    \STATE 加载 Canonical；要求目标高度早于执行高度，并匹配声明的
    基版本和前序摘要
    \STATE 检查条目顺序，以及输出和槽位的唯一性
    \STATE 对每个条目，要求声明位置存在决定 $\mathcal D_x$，
    且规范槽位保存 $x$ 的 \texttt{OriginalFunding}
    \STATE 重建获准的合并资金替换，并
    验证其开口
    \STATE 要求规范候选区块体的摘要等于
    $B.\mathsf{Next}$
    \STATE 验证非目标字节、长度，以及区块哈希和
    部件集头是否等于存储的原始身份
    \STATE $O\gets\operatorname{Overlay}(S)$；删除待处理条目，
    并记录逐输出引用
    \STATE 向 $O$ 添加一个逻辑修订版本、Task 和队列条目
    \STATE 返回该转移及不含经济效果的结果
  \end{algorithmic}
  \caption{表示批次的原子执行。失败会拒绝
    整个批次，且只有成功的转移才会提交。批次
    身份涵盖完整的规范命令，包括其证明
    数据。}
  \label{fig:supp-protocol-batch}
\end{figure}


\subsection{额外的历史读取开销}
对于 $(256,8)$ 第 1 轮运行中已解码的对象，两条路径上
单独重新检查授权的成本均约为每块 26.3\,ms，
输入 CH 检查为 10.3\,ms，完整 BlockID 重建为 0.67\,ms。
这些独立复查不计入可信读取方的计时。

\begin{figure*}[t]
  \centering
  \includegraphics[width=\textwidth]{reader-cost.pdf}
  \caption{读取方的权衡。(a,b) 三个运行级 P50 的中位数及
    最小值--最大值范围。(c) 主要逻辑记录的平均字节数。
    (d) 顺序执行各阶段的平均成本，而不是
    正文历史读取表中的总工作量中位数。}
  \label{fig:reader-cost}
\end{figure*}

\subsection{注入时延配置与逐轮运行结果}
\begin{table}[t]
  \caption{注入时延研究的逐轮运行结果（\texttt{9879f64}）。公开完成和成员完成的 P95 以 s 为单位；链耗时为该轮运行中两条 10 跳链的中位数和最大值，以 ms 为单位。}
  \label{tab:supp-faults}
  \centering
  \small
  \setlength{\tabcolsep}{3.5pt}
  \begin{tabular}{@{}lrrrrr@{}}
    \toprule
    运行 & \shortstack{接收\\P50/P95 (ms)} & \shortstack{公开完成\\P95} & \shortstack{成员完成\\P95} & \shortstack{链耗时\\中位数/最大值} & \shortstack{窗口内\\公开完成数} \\
    \midrule
    H-1   & 42.88/87.52 & 1.277  & 1.361  & 577/619  & -- \\
    H-2   & 41.56/83.20 & 1.288  & 1.375  & 599/619  & -- \\
    H-3   & 40.68/78.23 & 1.263  & 1.344  & 515/542  & -- \\
    L-1   & 81.91/118.22 & 1.315 & 1.410  & 946/980  & -- \\
    L-2   & 83.73/123.52 & 1.360 & 1.468  & 1001/1029 & -- \\
    L-3   & 78.05/111.80 & 1.307 & 1.394  & 815/835  & -- \\
    L-M-1 & 77.65/112.85 & 1.281 & 29.002 & 844/846  & 599 \\
    L-M-2 & 78.31/115.68 & 1.327 & 29.026 & 842/869  & 601 \\
    L-M-3 & 77.71/108.46 & 1.304 & 28.473 & 881/908  & 603 \\
    L-C-1 & 78.18/110.94 & 12.627 & 12.744 & 906/965 & 510 \\
    L-C-2 & 77.33/109.32 & 10.018 & 10.042 & 916/992 & 400 \\
    L-C-3 & 77.41/108.44 & 9.516  & 9.618  & 873/901 & 421 \\
    \bottomrule
  \end{tabular}
\end{table}
固定的四验证者配置运行一个组织的四名成员和一个网关。每个 TCP 方向使用一个有界 FIFO，最多容纳 64 个 16 KiB 数据块，并在到达时间加上所配置时延后释放。所有委员会对等节点地址均使用代理；当验证者保留出站连接时，对重复对等连接的抑制可能使监听代理处于空闲状态。钱包--网关通信、公开 HTTP、钱包交付和存储不注入时延。全部十二轮最终运行均使用 16,384 的文件描述符上限。一次预试运行继承了 256 的描述符上限后，重新执行了整个矩阵；这些受工具影响的预试结果单独保留，不参与合并统计。节点在离线数据库审计之前停止。

每轮运行中的两条链，其耗时均从首次发送计至最后一次认证接收。先在该轮内部取两条链的中位数，再对三次重复运行取中位数。故障窗口 QC 样本组中的请求在暂停至少一秒后发起，并在恢复之前完成接收。成员故障样本组中的全部 QC 恰好包含成员 0、1 和 2。离线验证检查委员会故障各轮运行中的全部 91、91 和 95 个已提交 BlockID。固定集合的提案者优先级表明，被暂停的验证者在每轮运行中有五次第 0 轮提案机会；对应高度均在第~1~轮提交，且不含该验证者的签名。窗口判定使用最晚的有效预提交时间戳，而不是前一高度的区块头时间。

\subsection{归档混合负载中从计划发送到接收的计时}
对于全部 84,000 笔普通付款，先计算记录的接收时间戳减去计划时间戳，再求分位数。沿用归档数据的最近秩计算约定，再取三个运行级 P95 的中位数，得到 N/R/C/P 分别为 304.253、140.007、475.236 和 666.915\,ms。实际发送至接收的 P95 分别为 152.766、139.894、136.491 和 144.205\,ms。前一种口径包含发送前的背压；两种口径都不是将单独的 P95 相加得到的。第 3 轮的部分墙上时钟记录显示发送早于计划时间戳（最小值为 $-13.104$\,ms）；原始差值和计数予以保留，不做截断。这些归档时间戳无法区分定时器唤醒误差与时钟调整，也不用于宣称亚毫秒级的发送调度精度。
